\section{Глава 17. Передовые рубежи}
\subsection{Упражнение 17.1(с. 526)}
В этом разделе описаны опции для случая с обесцениванием, но обесценивание вряд
ли применимо к управлению, когда используется аппроксимация функций
(раздел 10.4). Как выглядит естественное уравнение Беллмана для иерархической
стратегии, аналогичное (17.4), но в постановке со средним вознаграждением
(раздел 10.3)? Каковы две части модели опций, аналогичные (17.2) и (17.3), для
постановки со средним вознаграждением?
\begin{equation}
    r(s,\omega) \doteq \mathbb E\big[R_1+\gamma R_2+\gamma^2 R_3+\cdots+\gamma^{\tau-1}R_\tau
    \mid S_0=s,\ A_{0:\tau-1}\sim\pi_\omega,\ \tau\sim\gamma_\omega\big],
    \tag{17.2}
\end{equation}
\begin{equation}
    p(s'\mid s,\omega) \doteq \sum_{k=1}^{\infty}\gamma^k
    \Pr\{S_k=s',\ \tau=k\mid S_0=s,\ A_{0:k-1}\sim\pi_\omega,\ \tau\sim\gamma_\omega\},
    \tag{17.3}
\end{equation}
\begin{equation}
    v_\pi(s) = \sum_{\omega\in\Omega(s)}\pi(\omega\mid s)
    \Big[r(s,\omega)+\sum_{s'}p(s'\mid s,\omega)v_\pi(s')\Big],
    \tag{17.4}
\end{equation}
\begin{solution}
    \textbf{Части модели опций.} В постановке со средним вознаграждением
    обесценивание отсутствует, поэтому в (17.2) и (17.3) следует положить
    $\gamma=1$:
    \begin{equation}
        r(s,\omega) \doteq \mathbb E\big[R_1+R_2+\cdots+R_\tau
        \mid S_0=s,\ A_{0:\tau-1}\sim\pi_\omega,\ \tau\sim\gamma_\omega\big],
        \tag{17.2$'$}
    \end{equation}
    \begin{equation}
        p(s'\mid s,\omega) \doteq \sum_{k=1}^{\infty}
        \Pr\{S_k=s',\ \tau=k\mid S_0=s,\ A_{0:k-1}\sim\pi_\omega,\ \tau\sim\gamma_\omega\}.
        \tag{17.3$'$}
    \end{equation}
    Величина $r(s,\omega)$ --- полное вознаграждение, накопленное за время
    выполнения опции. В отличие от дисконтированного случая, (17.3$'$) снова
    является настоящим распределением: если опция завершается за конечное время
    почти наверное ($\Pr\{\tau<\infty\}=1$), то $\sum_{s'}p(s'\mid s,\omega)=1$,
    и обозначение «$\mid$» становится условным в обычном смысле.

    \textbf{Уравнение Беллмана.} В постановке со средним вознаграждением из
    вознаграждения на каждом \emph{элементарном} шаге вычитается $r(\pi)$ ---
    среднее вознаграждение иерархической стратегии $\pi$,
    \[
        r(\pi) \doteq \lim_{h\to\infty}\frac1h\sum_{t=1}^{h}
        \mathbb E\big[R_t\mid S_0,\ A_{0:t-1}\sim\pi\big],
    \]
    где усреднение ведётся по элементарным шагам, а не по моментам принятия
    решений. Опция занимает $\tau$ шагов, поэтому вычесть нужно не $r(\pi)$, а
    $r(\pi)$, умноженное на ожидаемую длительность опции:
    \begin{equation}
        \bar\tau(s,\omega) \doteq \mathbb E\big[\tau\mid S_0=s,\ A_{0:\tau-1}\sim\pi_\omega,\ \tau\sim\gamma_\omega\big].
        \tag{17.6}
    \end{equation}
    Дифференциальное уравнение Беллмана для иерархической стратегии принимает вид
    \begin{equation}
        v_\pi(s) = \sum_{\omega\in\Omega(s)}\pi(\omega\mid s)
        \Big[r(s,\omega)-r(\pi)\,\bar\tau(s,\omega)+\sum_{s'}p(s'\mid s,\omega)v_\pi(s')\Big],
        \tag{17.4$'$}
    \end{equation}
    где $v_\pi$ --- дифференциальные ценности состояний, определённые с точностью
    до общей аддитивной константы. При $\bar\tau\equiv 1$ (все опции ---
    низкоуровневые действия) уравнение сводится к обычному дифференциальному
    уравнению Беллмана.
\end{solution}
